% The same A_3 weave Delta Delta -> Delta as in
% a3_tetrahedral_weave.tex, now drawn as a top-to-bottom string diagram.
% The three four-valent commutations suppressed in the tetrahedral projection
% are drawn here as ordinary transverse crossings.
\begin{tikzpicture}[
  x=0.53cm,
  y=0.36cm,
  aThreeString/.style={line width=0.92pt,line cap=round,line join=round},
  aThreeStringOne/.style={aThreeString,draw=wqteal},
  aThreeStringTwo/.style={aThreeString,draw=wqorange,
    dash pattern=on 2.8pt off 1.35pt},
  aThreeStringThree/.style={aThreeString,draw=wqblue,
    dash pattern=on 0.72pt off 1.22pt},
  aThreeStringT/.style={circle,draw=black!82,fill=blue!12,
    minimum size=2.15mm,inner sep=0pt,line width=0.54pt},
  aThreeStringH/.style={circle,draw=black!82,fill=white,
    minimum size=3.05mm,inner sep=0pt,line width=0.64pt},
  aThreeStringVertexLabel/.style={font=\tiny,fill=white,
    inner xsep=0.55pt,inner ysep=0.25pt,text=black!82},
  aThreeStringBoundary/.style={font=\tiny,text=black!78},
  aThreeStringWord/.style={font=\scriptsize,text=black!84}
]
\path[use as bounding box] (-6.30,-1.70) rectangle (7.15,17.70);

% Boundary vertices for 121321 121321 and 121321.
\foreach \name/\x in {
  b0/-5.5,b1/-4.5,b2/-3.5,b3/-2.5,b4/-1.5,b5/-.5,
  b6/.5,b7/1.5,b8/2.5,b9/3.5,b10/4.5,b11/5.5}{
  \coordinate (\name) at (\x,16);}
\foreach \name/\x in {e0/-2.5,e1/-1.5,e2/-.5,e3/.5,e4/1.5,e5/2.5}{
  \coordinate (\name) at (\x,0);}

% The thirteen local moves, in sweep order.  Vertices in the same horizontal
% layer of the tetrahedral projection are aligned.  Each commuting crossing is
% placed at the actual intersection of its two straight strands.
\coordinate (v1)  at ( 0.0,15);
\coordinate (v2)  at ( 0.0,14);
\coordinate (v4)  at (-2.0,11.5);
\coordinate (v5)  at ( 1.0,11.5);
\coordinate (v7)  at ( 0.0, 8);
\coordinate (v8)  at (-2.0, 6.5);
\coordinate (v9)  at ( 2.0, 6.5);
\coordinate (v11) at (-3.5, 2);
\coordinate (v12) at (-0.35,2);
\coordinate (v13) at ( 2.5, 2);

% Exact straight-through locations of the three commuting crossings.
\coordinate (v3)  at ($(v2)!60/89!(v4)$);
\coordinate (v6)  at ($(v5)!1/6!(v9)$);
\coordinate (v10) at ($(v9)!305/421!(v12)$);

% Color 1 strands.
\foreach \u/\v in {
  b5/v1,b6/v1,v1/v2,b2/v4,v2/v3,v3/v4,v2/v5,b8/v5,
  v4/v8,v5/v6,v6/v9,b0/v11,v8/v11,v8/v12,
  v9/v10,v10/v12,v9/v13,b11/v13,
  v11/e0,v12/e2,v13/e5}{
  \draw[aThreeStringOne] (\u)--(\v);}

% Color 2 strands.
\foreach \u/\v in {
  b4/v2,b7/v2,v2/v7,v7/v8,v7/v9,b1/v8,b10/v9,
  v8/e1,v9/e4}{
  \draw[aThreeStringTwo] (\u)--(\v);}

% Color 3 strands.
\foreach \u/\v in {
  b3/v3,v3/v7,b9/v6,v6/v7,v7/v10,v10/e3}{
  \draw[aThreeStringThree] (\u)--(\v);}

% Six contractions and four A_2 braid vertices.
\foreach \v in {v1,v4,v5,v11,v12,v13}{\node[aThreeStringT] at (\v) {};}
\foreach \v in {v2,v7,v8,v9}{\node[aThreeStringH] at (\v) {};}

% Octahedron labels on the ten mutation vertices.
\node[aThreeStringVertexLabel,right=0.80mm of v1]  {$(0,0,2,0)$};
\node[aThreeStringVertexLabel,right=0.80mm of v2]  {$(0,0,1,1)$};
\node[aThreeStringVertexLabel,left=0.80mm of v4]   {$(1,0,1,0)$};
\node[aThreeStringVertexLabel,right=0.80mm of v5]  {$(0,1,1,0)$};
\node[aThreeStringVertexLabel,right=0.80mm of v7]  {$(0,0,0,2)$};
\node[aThreeStringVertexLabel,left=0.80mm of v8]   {$(1,0,0,1)$};
\node[aThreeStringVertexLabel,right=0.80mm of v9]  {$(0,1,0,1)$};
\node[aThreeStringVertexLabel,left=0.80mm of v11]  {$(2,0,0,0)$};
\node[aThreeStringVertexLabel,left=0.80mm of v12]  {$(1,1,0,0)$};
\node[aThreeStringVertexLabel,right=0.80mm of v13] {$(0,2,0,0)$};

% Boundary words and their individual colors.
\node[aThreeStringWord] at (0,17.35)
  {$\boldsymbol\Delta\boldsymbol\Delta$};
\foreach \x/\c in {
  -5.5/1,-4.5/2,-3.5/1,-2.5/3,-1.5/2,-.5/1,
  .5/1,1.5/2,2.5/1,3.5/3,4.5/2,5.5/1}{
  \node[aThreeStringBoundary,above=0.45mm] at (\x,16) {$\c$};}
\foreach \x/\c in {-2.5/1,-1.5/2,-.5/1,.5/3,1.5/2,2.5/1}{
  \node[aThreeStringBoundary,below=0.45mm] at (\x,0) {$\c$};}
\node[aThreeStringWord] at (0,-1.35) {$\boldsymbol\Delta$};
\end{tikzpicture}
